% Transcription — fonds Alexandre Grothendieck, shelfmark 85, batch 1 (pages 1-20).
% Established from the facsimile archives/batches/85/batch-01.pdf (a copy of
% 85.pdf fetched through the project's relay,
% https://grothendieck-relay.onrender.com/source/85.pdf; PDF page k is
% archivists' page k, no cover sheet), per the skill transcribe-grothendieck.
% The folder is 97 pages in five batches (1-20, 21-40, 41-60, 61-80, 81-97);
% this is the first, the others transcribed in parallel.
%
% Pass: Opus 5.5 (claude-opus-5-5), 2026-09-26 — first pass, unchecked against the pages by a human.
%
% Pages skipped: none. Pages 3 and 11 are wrappers inscribed in his hand and
% nothing else; by the cover rule they take no \page marker and their words
% become section headings, with a \note: p. 3 « Représentations projectives
% de G_2 (cas non dégénéré i.e. non diédral) », with « (7) » in pencil at the
% left; p. 11 « Représentations projectives de G_1 (et repr. projectives
% dégénérées de G_2) », with « (10 p) ». Both numbers are page counts of the
% runs that follow (4-10 is seven pages; 12-21 is ten, so the second run
% crosses into batch 2). Pages summarised: none; pages 12-20 are fast
% drafting under heavy cancellation, and their connective prose is largely
% \ill{}, following folder 29 batches 8-9.
%
% His own pagination: none on these leaves. The runs are followed by the
% wrappers' headings. On p. 1 his rubrics carry circled numerals, set as
% (1), (1'), (2), (4), (3) in the order they stand, with a \note. No dates
% in his hand.
%
% What the batch is about. P. 1 is a work plan (with « cf verso », i.e.
% p. 2): degenerate representations of G_0 = Z/2, of D_inf = G_1 and of G_2
% (via D_inf x Z/2 and (Z/2)^3), groups and maps attached to projective
% representations of G_2 over (finite) fields, genus-0 covers of the sphere
% minus three points, finite regular « hyperpolygônes », Pythagorean Schwarz
% systems, Gal(Qbar/Q). P. 2 (its verso, mostly illegible prose): any map
% of L_3 = {rho_0 rho_1 rho_2 = 1} to a finite group gives a finite map, so
% homomorphisms of L_2 into GP(1,k) give regular hyperpolygons. Pp. 4-10
% (run « (7) »): G_2 = <tau_0, tau_1, tau_2 | tau_i^2 = (tau_0 tau_2)^2 = 1>,
% rho_0 = tau_0 tau_1, rho_1 = tau_1 tau_2; a Theorem that the category C
% of (G, phi), G a form of GP(1)_S, phi : Gamma -> G(S) with rho_0s, rho_1s
% of square != 1 and not commuting, is rigid and (for Gamma = G_2)
% equivalent to the discrete category of pairs (alpha_0, alpha_1) of
% sections with alpha_0 + 2, alpha_1 + 2, alpha_0 + alpha_1 invertible,
% alpha_i = A(rho_i) the invariant; proof via the regular 2-polyhedron,
% 4cos^2 theta_i = beta_i, alpha_i = 2cos 2theta_i, and a lemma
% A(rho_Y) = 2cos theta for a rotation of a rank-3 quadratic module; a
% Corollary for Gamma_{nu_0,nu_1} (orders exactly nu_i, Psi_{nu_i}(alpha_i)
% = 0, a table of three cases on invertibility of nu_i); a cancelled
% remark on nu_i = 2p; examples with Gamma finite: A) (3,3), S_4, empty
% in char. 2, unique over Z[1/2]; B) (4,3), Z/2 x S_4; C) (3,5),
% Z/2 x A_5, alpha^2 + alpha - 1 = 0, a Theorem. Pp. 12-20 (run
% « (10 p) », continuing into batch 2): Prop. on tau_i in GP(1,k), k alg.
% closed — A) a stable degree-2 divisor, A') invariant line in the conic
% plane C, B) tau_i in N(T) or B, C'') factorisation through an extension
% of cyclic groups, C') through D_inf x Z/2, D) two of the three pairs
% commute — with a sketch of the implications; a Lemma on the subgroups
% N_i, N_{i,j} of D_inf and D_i, Cor. 1-2 (quotients of D_inf and of
% D_inf x Z/2; D_i x Z/2 = D_2i for i odd); a Proposition on faithful
% D_i -> GP(1,k) (p = 0 always; i = 0 iff k^* has an element of infinite
% order; i >= 2 iff p does not divide i or i = p); a cancelled Cor. on
% D_i x Z/2; then representations of D_inf = G_1 in forms of GP(1)_S with
% u_s^2 != 1: a Theorem that C is equivalent to the category C' of pairs
% (alpha, R), alpha + 2 invertible, R a torsor under U(2 - alpha), where
% U(rho) is the group of lambda with 1 + lambda rho invertible under
% [lambda][lambda'] = [lambda + lambda' + lambda lambda' rho]; the quasi-
% inverse via H_1 \ {c u d}; Aut(G, phi) = 2(U_rho); Cor. 1 (rho
% invertible: U_rho = G_m, image the maximal torus through u, Aut = mu_2),
% Cor. 2 (alpha = 2: U_rho = G_a, Aut trivial if 2 invertible, = G_a in
% char. 2), Cor. 3 for Gamma = D_inf x gamma (triples (alpha, R, phi_0));
% a Proposition on D_i x Z/2 -> GP(1,k), and the conditions for phi to
% factor through a mono D_n -> G: n invertible or n = p odd prime, and
% Psi_n(alpha) = 0, Psi_n a « ½cyclotomique » polynomial.
%
% Notation fixed once for the batch, following folders 76, 82 and 88:
% GP(1) as \mathrm{GP}(1), his underlined Aut as \underline{\mathrm{Aut}};
% the dihedral groups in his double-struck D as \mathbb{D}_n,
% \mathbb{D}_\infty; the cartographic groups G_1, G_2 as he writes them;
% his circled S and A as \mathfrak{S}_4, \mathfrak{A}_5; G_m, G_a as
% \mathbb{G}_m, \mathbb{G}_a; the invariant map A : G -> E^1_S as written
% (the E is a doubtful reading); the centre on p. 9 as \mathfrak{z};
% {}_2 gamma, {}_2 U for 2-torsion, the 2 written low at the left.
%
% Hardest pages: 2 (prose almost wholly lost), 12-16 and 18 (layered
% cancellations), and the left-margin notes of 17 and 18. Readings to check
% first: « E^1_S » on p. 4; « f_i f_{i+1} » in the NB of p. 5; the
% whole first paragraph of p. 6 and the factor (alpha_0 - 2) there; the
% relations in C', C'' and the proof sketch on pp. 12-13; the conditions
% of the Proposition on p. 15; « zAnt » on p. 17; p. 18 throughout.

\documentclass[11pt,a4paper]{article}
\input{../preamble/grothendieck}

\folder{85}
\batch{1}
\pages{1}{20}
\dating{s.d. — le groupe « Géométrie et topologie combinatoire » (69 à 102) est daté 1976-[vers 1986]}
\watermark{Édition de démonstration}
\foldertitle{2-polyèdres et 1-polygônes : notes manuscrites (s.d.)}

\begin{document}

\page{1}

\emph{Opérations} du centre de $\Gamma = G_2/-$
\note{la page est un plan de travail ; ses rubriques sont numérotées par des
chiffres entourés, rendus ici (1), (1'), (2), (4), (3), dans l'ordre où elles
se suivent sur la page. Chaque rubrique est reliée par une accolade à son
numéro, et séparée de la suivante par un trait}

(1)\quad Représentations \add{n.\ dég.} \struck{\ill{}} \add{\uncertain{singulières}}
\quad $G_0 = \mathbb{Z}/2$

Représentations \emph{dégénérées} de $\mathbb{D}_\infty = G_1$ --- i.e.\ repr.\
\add{fidèles} de $\mathbb{D}_2$ / relations \ill{} : un octaèdre, un
\uncertain{quaternions} / quaternions généralisés\ldots.

Représentations \uncertain{simplement} dégénérées de $G_2$ --- i.e.\
représentations \uncertain{non dég.}\ de $\mathbb{D}_\infty \times \mathbb{Z}/2$
(\struck{considérées} \add{\uncertain{d'après un}} \ill{} les cartes
correspondantes, et leur genre\ldots)

Représentations \uncertain{doublement} dégénérées de $G_2$ --- i.e.\ repr.\ de
$\mathbb{D}_2 \times \mathbb{Z}/2 \simeq (\mathbb{Z}/2)^3$ avec des « rang »
\struck{\ill{}} $\geqslant 2$

(1')\quad Pts de ramification des repr.\ \struck{de} $G_2$\ldots
\note{la rubrique (1') est marquée dans la marge d'un triple trait vertical}

Groupes et cartes associés aux représentations \add{projectives} de $G_2$ sur
un corps --- notamment sur un corps fini

Les rev.\ \struck{de degré} \add{et les \uncertain{modèles}} de genre 0 de la
sphère privée de trois points

Les cartes finies associées aux hyper-polygônes \emph{réguliers} finis (cf
verso)
\marginal{cf verso}
\note{« cf verso » est écrit une seconde fois, verticalement, dans la marge
gauche, devant l'accolade qui réunit ces trois alinéas ; le verso est la
page 2}

(2)\quad Terminer en regardant les syst.\ de Schwarz pythagoriciens (cas de
$G_5$), et déterminer \struck{l'explication} la structure d'ens.\ : gr.\
sous les \uncertain{gps} des gr.\ \uncertain{riemanniens}, dans les trois cas

(4)\quad structure des \ill{} des gr.\ \uncertain{riemanniens} et relations,
\ill{} $\mathrm{Gal}(\overline{\mathbb{Q}}/\mathbb{Q})$

(3)\quad Détermination des hyperpolygônes \add{pythagoriciens}
(\ill{} l'ambiguïté \ill{}\ldots)
\note{la rubrique (3) est écrite en dernier, tout au bas de la page, sous un
trait, après (4)}

\page{2}

Toute \struck{\uncertain{représentation}} \add{\ill{}} du \uncertain{gpe}
\ill{}
\[
L_3 = \lbrace \rho_0, \rho_1, \rho_2 \mid \rho_0\rho_1\rho_2 = 1 \rbrace
\]
dans un gpe fini $\Gamma$ donne une \uncertain{carte} finie
\struck{\ill{}} 3-\ill{} \ill{} (\ill{} \uncertain{régulière},
\uncertain{génériquement}) --- \ill{} \ill{} \ill{}, \ill{} déf.\ Donc tt
\uncertain{homom} de $L_2$ dans un $\mathrm{GP}(1, k)$, \struck{\ill{}} \ill{}
\ill{} fini, donne une telle \ill{} --- \ill{} les représentations en question
\ill{} une \uncertain{interprétation} géométrique \ill{} en termes
d'hyperpolygônes \emph{réguliers}. Donc on a une étude
\uncertain{cristallographique} de ceux-ci --- un \ill{} \ill{} par ceux qui
les définissent par des systèmes de Schwarz \struck{type} pythagoriciens.
Ceux qui définissent \ill{} 44 types d'« hyperpolygônes » --- obtenus par
\struck{\uncertain{symétries}} et « \uncertain{quintiques} » \ill{} (plus
petits c'est fort gros !), \ill{} \ill{} \struck{\uncertain{gars}}
(\uncertain{gr.} \struck{\ill{}} fini de \ill{} \ill{} tous.
\struck{D'autres ex.\ \ill{}} \ill{} ceux-ci \ill{} \ill{}
\uncertain{correspondants} \ill{} \ill{} éléments \emph{canoniques} des
$n$-\ill{} \uncertain{hyperboliques} de $\mathrm{GP}(1)$ \ill{}, aux polyèdres
réguliers \ill{} --- plus leurs transformés par \ill{}.
\note{l'indice 3 de $L_3$ surcharge un 2 ; plus bas il écrit $L_2$. La page
est d'une écriture rapide dont les formules seules passent sûrement ; la
prose est en grande partie perdue. « cartes » (première phrase), « Schwarz »,
« pythagoriciens » et « hyperpolygônes » sont confirmés par le plan de la
page 1}

\marginal{Il convient d'y ajouter ceux qui sont « dégénérés », i.e.\ qui
correspondent \ill{} trois \ill{} \ill{} aux représentations de
$\mathbb{D}_\infty \times \mathbb{Z}/2$, et types dégénérés (représentations
\ill{} de $G_2$) et \ill{} \ill{} \ill{} \ldots}
\note{écrit verticalement dans la marge gauche, sur quatre lignes}

\section*{Représentations projectives de $G_2$ (cas non dégénéré i.e.\ non
diédral)}

\note{inscrit en haut à droite du feuillet de garde (page 3), qui ne porte
rien d'autre et ne reçoit pas de numéro de page. À gauche, au crayon, « (7) »
: sans doute le nombre de feuillets de la liasse, qui compte en effet sept
pages de texte (pages 4 à 10)}

\page{4}

Soient $\Gamma$ un groupe, $(\tau_i)_{0 \leqslant i \leqslant 2}$
\struck{\ill{}} syst.\ de générateurs de $\Gamma$, avec $\tau_i^2 = 1$,
$\tau_0\tau_2 = \tau_2\tau_0$ \struck{\ill{}}. On pose $\rho_0 = \tau_0\tau_1$,
$\rho_1 = $ \struck{\ill{}} $\tau_1\tau_2$.

\emph{Thm} $S$ schéma, \struck{\ill{} $\mathrm{GP}(1)_S$} \struck{\ill{}}
\struck{\ill{}} Considérons :
\note{le premier énoncé, après « $S$ schéma », est biffé sur toute la ligne ;
au-dessus, « \uncertain{Soit} » ajouté puis biffé lui aussi}

A) Catégorie $C$ des couples $(G, \varphi)$, où $G$ est une forme de
$\underline{\mathrm{GP}}(1)_S$, $\varphi : \Gamma \to G(S)$, satisfaisant les
conditions

(*) $\forall s \in S$, $\rho_{0s}^2$ et $\rho_{1s}^2$ sont $\neq 1_s$
[NB on désigne par $\tau_0, \tau_1, \tau_2, \rho_0, \rho_1$ les images de
$\tau_0, \tau_1, \tau_2, \rho_0, \rho_1$ dans $G(S)$] et $\rho_{0s}$ et
$\rho_{1s}$ ne commutent pas

\struck{($\Leftrightarrow$ b) $\forall s \in S$, l'image $\Gamma_s$ de $\Gamma$
dans $G(s)$ [$\simeq \Gamma/N_s$ où $N_s$ \ill{} \ill{} de $\Gamma \to G(s)$]
n'est pas isomorphe à un groupe diédral (c.-à-d.\ un produit
$\mathbb{D}_\infty \times \mathbb{Z}/2$ \ill{}, $\mathbb{D}_n$ ($n \geqslant 3$)
--- le cas $n = 1, 2$ est déjà exclu par (*), qui signifie \ill{} que
$\tau_{0s}$ et $\tau_{1s}$ \ill{}, ni $\tau_{1s}$ et $\tau_{2s}$ --- donc
$\Gamma_s$ \ill{} \ill{} \uncertain{abélien}.)}
\note{tout ce passage b) est barré de grandes croix ; une partie en est en
outre biffée ligne à ligne (« un produit $\mathbb{D}_\infty \times
\mathbb{Z}/2$ \ldots », « $\mathbb{D}_n$ ($n \geqslant 3$) »)}

\marginal{\struck{On peut remplacer b) par b') $\rho_0$ et $\rho_1$ ne
commutent pas (\uncertain{ou encore} $\rho_0$ et $\sigma = \rho_0\rho_1 =
\tau_0\tau_2$ ne commutent pas) (NB $\sigma^2 = 1$)}}
\note{dans la marge gauche, à la hauteur de b) ; barré avec lui}

Les flèches de $C$ sont les isomorphismes.

B) La catégorie discrète $C'$ définie par l'ens.\ des couples
$(\alpha_0, \alpha_1) \in \Gamma(S, \underline{\mathcal{O}}_S)^2$,
satisfaisant les conditions

(*) \struck{\ill{}} $\alpha_0 + 2$, $\alpha_1 + 2$, $\alpha_0 + \alpha_1$ sont
inversibles.
\note{sous (*), une première rédaction de la condition, surchargée, est
biffée ; le couple $(\alpha_0, \alpha_1)$ de B) est surchargé, et semble
avoir été écrit d'abord $(\alpha_0(\varphi), \alpha_1(\varphi))$}

Pour tout objet $(G, \varphi)$ de $C$, soit
\[
\alpha_0 = A(\rho_0), \qquad \alpha_1 = A(\rho_1)
\]
où
\[
A : G \longrightarrow \mathbb{E}^1_S
\]
est l'invariant. \struck{On trouve} Ceci posé :

1) $(\alpha_0, \alpha_1)$ ainsi associé : $(G, \varphi)$ satisfait (*),

2) Le foncteur
\[
C \longrightarrow C'
\]
défini ainsi par $(G, \varphi) \longmapsto (\alpha_0(\varphi), \alpha_1(\varphi))$
est \emph{pleinement fidèle}, et c'est une équivalence de catégories
\add{si $\Gamma = \lbrace \tau_0, \tau_1, \tau_2 \mid \tau_0^2 = \tau_1^2 =
\tau_2^2 = (\tau_0\tau_2)^2 = 1 \rbrace = G_2$},
\struck{\ill{} et autres \ill{}}
\note{la ligne « si $\Gamma = \ldots = G_2$ » est écrite au-dessus de la ligne,
à la place de trois mots biffés}

a) La catégorie $C$ est \emph{rigide}, i.e.\ $\forall (G, \varphi)$, le
centralisateur de $\varphi$ dans $G(S)$ est $\lbrace 1 \rbrace$.

b) Deux objets $(G, \varphi)$ et $(G', \varphi')$ de $C$ sont isomorphes ssi
$\alpha_0(\varphi) = \alpha_0(\varphi')$, $\alpha_1(\varphi) = \alpha_1(\varphi')$

3) \struck{$\forall$} \add{si $G = G_2$} \uncertain{couple} $(\alpha_0, \alpha_1)$
\struck{\ill{}} de sections de $\underline{\mathcal{O}}_S$ satisfaisant (*),
$\exists (G, \varphi)$ \struck{\ill{}} objet de $C$ tel que
$\alpha_0(\varphi) = \alpha_0$, $\alpha_1 = \alpha_1(\varphi)$ [et ledit objet
de $C$ est défini à isom.\ unique près grâce à a)]
\note{au début de 3), un « $\forall$ » est suivi d'une courte insertion
au-dessus de la ligne, « si $G = G_2$ » (lecture incertaine : sans doute
$\Gamma = G_2$), enfermée dans un trait qui la biffe peut-être aussi}

\page{5}

\emph{Démonstration}\quad 1) On sait, par les généralités, que les conditions
a) b) impliquent \add{\uncertain{parmi autres}} l'existence d'un 2-polyèdre
\emph{régulier} épinglé, donnant lieu à la repr.\
\[
\hat{\varphi} : G_2 \ (\text{groupe cartographique}) \longrightarrow \Gamma
\xrightarrow{\ \varphi\ } G(S),
\]
\struck{d'où on déduit, on l'a vu, que $\alpha_0 + 2$, $\alpha_1 + 2$,
$\alpha_0 + \alpha_1$ sont inversibles.} défini à isom.\ unique près. Un
calcul immédiat donne que les invariants $\alpha_0, \alpha_1$ de ce
2-polyèdre sont les mêmes que $\alpha_0(\varphi)$ et $\alpha_1(\varphi)$. En
effet, dans le contexte général (dim $n$) avec paramètres $\beta_i = \alpha_i
+ 2$, pour $0 \leqslant i \leqslant n-1$,
\[
f(a_i) = (\beta_0 \cdots \beta_{i-1}), \quad f(a_{i+1}) = \beta_0 \cdots
\beta_i, \quad \varphi(a_i, a_{i+1}) = -(\beta_0 \cdots \beta_i)
\]
donc si les $\beta_i$ sont \struck{inv.} \ill{} (on peut \struck{\ill{}}
parler de « l'angle » $\theta_i$ \add{ainsi que \ill{}}) de $a_i$ et
$a_{i+1}$, \ill{} avec
\[
4\cos^2\theta_i = \frac{\varphi(a_i, a_{i+1})^2}{f(a_i)f(a_{i+1})} = \beta_i
\qquad \Bigl(\text{\emph{NB}}\ \xi_i = \frac{\beta_i}{f_i f_{i+1}} =
\frac{\beta_i}{4} = \cos^2\theta_i\Bigr)
\]
i.e.
\[
\alpha_i = 2(\underbrace{2\cos^2\theta_i - 1}_{\cos 2\theta_i}) = 2\cos 2\theta_i
\]
\note{la parenthèse NB est lue telle qu'elle se présente ; le dénominateur
« $f_i f_{i+1}$ » et l'égalité qui suit sont d'une lecture incertaine}

Il faut aussi prouver ceci :

\emph{Lemme} (\struck{\ill{}}) Soit $(V, f)$ un module quadratique libre de
rang 3, considérons une rotation \struck{\ill{}} $\rho$ d'angle
\struck{\ill{}} $\vartheta$ (autour d'un axe totalement non isotrope)
\ill{} soit $\alpha = 2\cos\vartheta$. \struck{Alors} Soit \struck{\ill{}}
$Y \subset \check{\mathbb{P}}(V)$ \uncertain{fibré} des droites projectives
définies par $f$, et considérons l'automorphisme $\rho_Y$ induit par $\rho$.
Ceci posé, l'invariant de $\rho_Y$ est donné par
\[
A(\rho_Y) = \alpha
\]
\note{dans la marge gauche, un petit croquis : deux droites se coupant, l'angle
$\theta_i$ marqué entre elles, et un arc fléché marqué $\rho_i$ (la rotation) ;
non redessiné. Au début de la ligne « Soit \ldots », « Alors » est biffé ;
avant « d'angle $\vartheta$ » un mot surchargé est biffé}

\emph{Dém} Quitte à choisir des axes de coordonnées convenables
\struck{\ill{}}, on peut supposer
\[
f(x, y, z) = x^2 + yz, \qquad \rho_\xi(x, y, z) = (x, \xi y, \xi^{-1}z),
\qquad \text{et } \alpha = \xi + \xi^{-1}.
\]
Or l'homom.\ $\mathbb{G}_m \hookrightarrow G = \mathrm{Aut}\, C$, $\xi \mapsto
\rho_\xi$, est \ill{}, donc on trouve $A(\rho_\xi) = \xi + \xi^{-1} = \alpha$,
cqfd.

\page{6}

On a un peu salopé la présentation, qui \uncertain{omet} \ill{} pour la
\ill{} \ill{} cas.\ \struck{\ill{}} $\neq 2$ (sans parler d'une « théorie des
angles » sur schémas quelconques, que je n'ai pas bien explicitée). Mais il
est clair \uncertain{par principe} de prolongement que si le formalisme
\add{\uncertain{sur l'infini}} est prouvé en car.\ 0 (bien \ill{} on se sert
\uncertain{sur terrain} \ill{}) il l'est en car.\ quelconque \ldots

Comme on est dans les cas de forme \add{\uncertain{quand}} \uncertain{invariante}
(i.d.\ \uncertain{réalisable} sur \struck{\ill{}} \add{\uncertain{anneau} fixé})
lisse, on doit bien avoir
$(\alpha_0 + 2)(\alpha_0 - 2)(\alpha_0 + \alpha_1)$ inversible.
Inversement, si on part d'un couple $(\alpha_0, \alpha_1)$ satisfaisant à ces
conditions, on sait qu'il \uncertain{associe} un polyèdre régulier,
\uncertain{d'où} une \uncertain{forme lisse}, d'où représentation de $G_2$,
correspondant : \ill{} \ill{}. \struck{On sait que la cat.}
\note{le produit se lit bien « $(\alpha_0 + 2)(\alpha_0 - 2)(\alpha_0 +
\alpha_1)$ » ; les pages 4 et 7 écrivent $(\alpha_0 + 2)(\alpha_1 + 2)(\alpha_0
+ \alpha_1)$}

\struck{A) \ill{}}

\emph{Cor.} \ill{}\quad Soient $\nu_0, \nu_1$ deux entiers $\geqslant 3$
\add{ou infinis}, et supposons que $\Gamma$ soit le groupe \uncertain{complet}
de relations
\[
\tau_0^2 = \tau_1^2 = \tau_2^2 = (\tau_0\tau_2)^2 = (\tau_0\tau_1)^{\nu_0} =
(\tau_1\tau_2)^{\nu_1} = 1
\]
i.e.\ $\Gamma = \Gamma_{\nu_0, \nu_1}$. \struck{Alors l'image \ill{} du
foncteur} Soit \struck{\ill{}} $C(\nu_0, \nu_1)$ la sous-catégorie de $C$
formée des $(G, \varphi)$ ($G$ forme de $\mathrm{GP}(1)_S$, $\varphi : \Gamma
\to G(S)$ \ill{}) telles que l'on ait
\struck{$\rho_0^{\nu_0} = \rho_1^{\nu_1} = 1$, \ill{}}

a) $\forall s \in S$, $\rho_{0s}$ est d'ordre $\nu_0$, $\rho_{1s}$ d'ordre
$\nu_1$ \emph{exactement} \struck{$\rho_0$ et $\rho_1$ ne commutent pas
(\ill{})}

\struck{b) $\forall s \in S$, $\Gamma_s$ n'est pas \ill{} à un $\mathbb{D}_n$
(\ill{} $\mathbb{D}_\infty \times \mathbb{Z}/2$)}
\struck{[\emph{NB} Je reviens au même d'exiger que $\rho_{0s}$ et $\rho_{1s}$
ne commutent pas, \ill{} il suffit pour \ill{}, compte tenu de a), que l'on
ait : on n'a pas $\nu_0$ impair et $\nu_1 = 2\nu_0$ ou l'inverse]}
\note{b) et la remarque entre crochets sont barrés de traits obliques}

\marginal{\ill{} que, $i = 0, 1$, $\nu_i$ \ill{} \ill{} si $\nu_i$ premier
\ill{} ($\nu_i \geqslant 3$)}
\note{note de la marge gauche, encadrée, reliée par un trait à « Soit » ;
lecture très incertaine}

\page{7}

Soit $C'(\nu_0, \nu_1)$ la cat.\ discrète définie par les couples $(\alpha_0,
\alpha_1)$, \struck{satisfaisant chaque que $\alpha_i$ satisfait} \ill{}
\struck{\ill{}}

a) $(\alpha_0 + 2)(\alpha_1 + 2)(\alpha_0 + \alpha_1)$ inversible

et chaque $\alpha_i$ satisfaisant les conditions suivantes

b) $\Psi_{\nu_i}(\alpha_i) = 0$ \qquad $i = 0, 1$

\struck{Supposons enfin que, \ill{} l'un des \ill{} l'un des $\nu_i$ est tel
que $\nu_i = 2p$ (où $p$ est premier impair)}

\struck{c) Si $\nu_i$ \ill{} de la forme $2p$ ($p$ premier) \ill{} \ill{} $p$
(\ill{} premier) est inversible, ou de la forme $2p$ avec $2$ \uncertain{inv.}
si $p$ impair, ou $\nu_i = p$ ($p$ premier impair).}
\note{c) est barré de traits obliques ; il avait été récrit plusieurs fois,
avec insertions et ratures, avant d'être abandonné. Dans la marge gauche, en
face, une note elle aussi barrée : « \ill{} ce qui est \ill{} \ill{} sur
$\nu_0, \nu_1$ \ill{} \ill{} »}

Alors on a une équivalence de catégories
\[
C(\nu_0, \nu_1) \longrightarrow C'(\nu_0, \nu_1).
\]


En effet, il faut \uncertain{remarquer}, \struck{\ill{}} \add{si $\Gamma =
\mathcal{G}$}, \ill{} $(G, \varphi)$ \add{\uncertain{$\in C$}} et $(\alpha_0,
\alpha_1)$ se correspondent, que \struck{A) les conditions a) b) c) sont
satisfaites, \ill{}} alors on a la condition \struck{a) \ill{} sur $\rho_0$ et
$\rho_1$}

(*) \quad $\rho_i^{\nu_i} = 1$ \quad et \quad $\forall s \in S$, $\rho_{is}$
d'ordre $\nu_i$ exactement

\struck{\ill{}} ssi ---

(**) \quad $\Psi_{\nu_i}(\alpha_i) = 0$ \struck{et c) \ill{} $\nu_i$}

\marginal{\struck{réf}}
\struck{Or il est} Or on sait que \add{(pour $\nu_i \geqslant 3$)}
\struck{pourvu que $\nu_i$ \ill{} de la forme $2p$ ($p$ premier) ou $p$ ($p$
premier impair), l'un ou l'autre,} que \ill{} s'exprime par les conditions
$\Psi_{\nu_i}(\alpha_i) = 0$ sur l'invariant $\alpha_i$ de $\rho_i$, plus la
condition $\nu_i$ \uncertain{inv.}\ si $\nu_i$ non premier :

\struck{(cas 1°) 2 inv.\ sur $S$ \ill{} \ill{}} \add{\ill{} $\nu_i$ inv.}

(cas 2°) aucune condition \ill{}.

(cas 3°) $\nu_i$ inv.\ sur $S$ \quad (\ill{} \ill{})

Ce qui prouve l'équivalence.
\note{les trois cas sont disposés en tableau, les numéros entre parenthèses
reliés par une accolade ; la première ligne est biffée et remplacée au-dessus
par la condition « $\nu_i$ inv. »}

\page{8}

\struck{\emph{Remarque} Comme \ill{} $\alpha_i + 2$ est inv.\ i.e.\ $\alpha_{i,s}
\neq -2$ en \ill{} $s$, \ill{} indique qu'en fait, si \ill{} $\nu_i = 2p$
\ill{} \ill{}\ldots}
\note{une première suite de la remarque, sur trois lignes, est biffée ligne à
ligne : « Alors \ill{} \ill{} \ill{}, \ill{} \ill{} \ill{} \ill{} et si $S$
contient un pt $s$ de \ill{} $p$ », puis « Alors $C'(\nu_0, \nu_1)$ est
\emph{vide} »}
\struck{On voit directement \ill{} dans ce cas $C(\nu_0, \nu_1)$ est
\emph{également} vide : en effet, un gr.\ diédral de $G_s$ (forme de
$\mathrm{GP}(1)_{k(s)}$) ne peut jamais \ill{} d'ordre $2p$. Donc il est
\add{\uncertain{raisonnable}}, \ill{} assurer que $C(\nu_0, \nu_1)
\xrightarrow{\ \sim\ } C'(\nu_0, \nu_1)$ soient \ill{}, d'exiger que $\nu_i =
2p$ \add{($p$ pr.\ impair)} implique que $p$ \ill{} sur $S$, donc \ill{} fait
$2p = \nu_i$ \ldots\ c'est-à-dire qu'on exige que $\nu_i$ est inv.\ sur $S$
\ill{} \ill{} $p$ \emph{premier} (impair ou $\nu_i \geqslant 3$).}
\note{toute la remarque est barrée de deux longs traits obliques. En marge,
en bas à gauche, un mot biffé, \uncertain{« faux »}, relié par un trait
vertical au bas de la page}

Exemples avec $\Gamma_{\nu_0, \nu_1}$ \emph{fini}

A) $\nu_0 = 3$, $\nu_1 = 3$, $\Gamma_{\nu_0, \nu_1} \simeq \mathfrak{S}_4$. On
\struck{trouve que} \struck{la catégorie des repr.\ de $\mathfrak{S}_4$
\ill{}} \ill{} restriction à faire sur $S$. On trouve que la cat.\ des
couples $(G, \varphi)$, $G$ forme de $\mathrm{GP}(1)_S$,
$\varphi : \mathfrak{S}_4 \to G(S)$ telles que $\forall s \in S$, $\rho_{0s}$
et $\rho_{1s}$ soient d'ordre 3 exactement \add{et ne commutent pas}, est
\add{\uncertain{vide si}} \struck{équivalente \ill{}} 2 (${} = (-1)(-1)(-1-1)$
au signe près) n'est pas inversible, et qu'elle est équivalente à la cat.\
ponctuelle dans le cas contraire. C'est-à-dire : si \ldots
\note{il note $\mathfrak{S}_4$ d'un S cerclé, sa graphie du groupe
symétrique ; le produit $-2 = (-1)(-1)(-1-1)$ est la valeur de $(\alpha_0 +
2)(\alpha_1 + 2)(\alpha_0 + \alpha_1)$ pour $\alpha_0 = \alpha_1 = -1$}

a) \struck{\ill{}} \add{sur un \ill{} de} car.\ 2, \uncertain{toute}
\uncertain{repr.}\ $\mathfrak{S}_4 \to G(S)$, dans \ldots\ \struck{\ill{}
$\rho_0$ et $\rho_1$ commutent, i.e.\ $\rho_0 \ldots \rho_1 = 1$ (i.e.\
$\tau_0 = \tau_1$ ou $\tau_1 = \tau_2$ \ill{})} \add{la} repr.\ se factorise
par \struck{$\mathfrak{S}_3$ ($= \mathbb{D}_3$)} \struck{\ill{}}
$\mathfrak{S}_3$ \struck{\ill{}}
\note{fin de page très raturée : plusieurs symboles $\mathfrak{S}_3$,
$\mathbb{D}_3$, $\mathfrak{S}_4$ sont biffés ou surchargés ; seul le second
$\mathfrak{S}_3$ paraît maintenu}

\page{9}

b) Sur un \struck{corps de car.} \add{schéma $S$ où 2 est inv.} (i.e.\ sur
$\mathbb{Z}[\frac{1}{2}]$), il y a à isom.\ unique près, un couple $(G,
\varphi)$ et un seul, où $G$ est une forme de $\mathrm{GP}(1)_S$, $\varphi :
\mathfrak{S}_4 \to G(S)$, avec $\rho_{0s}$ et $\rho_{1s}$ \struck{d'ordre 3
exact.}\ \add{\ill{} \ill{}} \struck{i.e.} \add{i.e.\ la repr.\ ne se factorise
pas par $\mathfrak{S}_3$} en chaque fibre. \struck{\ill{}} \add{\ill{}}
\uncertain{faut}, \ill{} \ill{} que ladite représentation est
\struck{\ill{}} \uncertain{partout} \uncertain{fidèle}]
\note{le crochet fermant répond peut-être à un crochet ouvert à la fin de la
page 8 ; les insertions au-dessus de la ligne sont serrées et en partie
illisibles}

B) $\nu_0 = 4$, $\nu_1 = 3$ (ou l'inverse) --- \struck{\ill{} \ill{}}
\add{ici \uncertain{on sait} que}
$\Gamma = \mathfrak{S}_4 \times \mathbb{Z}/2$. Le \uncertain{résultat}
\uncertain{général} nous dit que \add{\uncertain{inversible}} (\ill{} car.\ 2)
qu'il faut \uncertain{écarter}. [On pourrait déduire un \uncertain{énoncé}
\ill{} concernant le gpe $\Gamma$ de A), mais en passant la \ill{}
\uncertain{diédrale}, \ill{} \ill{} des renseignements \ill{}
\marginal{\ldots\ vraiment !}
\note{« vraiment ! » est écrit en biais dans la marge gauche, souligné, à
hauteur de cette ligne}
Manquent $\tau_0, \tau_1, \tau_2$ \ill{} \struck{\ill{}} \add{les}
gpes \emph{tétraédraux} (\uncertain{un plus} tétraédraux) \ill{}
$\rho_0^4 = \rho_1^3 = 1$, et \ldots\ \uncertain{forme}

a) Sur un corps de car.\ 2, \struck{$\forall$} hom.\ \add{$\varphi : \Gamma
\to G$} de $\Gamma \simeq \mathbb{Z}/2 \times \mathfrak{S}_4$ dans une forme de
$\mathrm{GP}(1)_k$,

\struck{\ill{} $\rho_0^2 = 1$ \ill{} l'ordre de $\rho_0$ est égal à 2 \ill{}
\ill{} $\rho_1 = \mathrm{id}$ \ill{} $\rho_0$ et $\rho_1$ commutent (i.e.\
\ill{} $\tau_0$ et $\tau_1$ commutent donc $\tau_2$ : \ill{} \ill{}
factorise par $\mathbb{Z}/2 \times \mathfrak{S}_3$, \ill{} si \ill{} $\tau_1 =
\tau_2$ \ill{} \ill{} \ill{} factorise par $\mathbb{Z}/2 \times
\mathfrak{S}_4/\mathbb{Z}_2$}
\note{tout ce bloc est barré de traits obliques et entouré d'une accolade ;
dans la marge gauche, reliée à l'accolade : « i.e.\ \ill{} \uncertain{se}
factorise par $\mathbb{Z}/2 \times \mathfrak{S}_3 \simeq \mathbb{D}_6$ »}

b) Sur un schéma $S$ où $2 \cdot 1_S$ est inv.,

$\exists !$ à isom.\ unique près couple $(G, \varphi)$, \add{telle que
$\varphi_s$ ne se factorise pas par $\mathbb{Z}/2 \times \mathfrak{S}_3 \simeq
\mathbb{D}_6$} où

$\varphi : \Gamma \to G$ forme de $\mathrm{GP}(1)$ --- et en fait

$\varphi$ est trivial sur $\mathbb{Z}/2 = \mathfrak{z}(\Gamma)$, et induit un
mono $\mathfrak{S}_4 \hookrightarrow G_s(S)$ $\forall s \in S$.
\note{le centre de $\Gamma$ est noté d'une petite lettre ronde, rendue
$\mathfrak{z}$}

\page{10}

C) $\nu_0 = 3$, $\nu_1 = 5$, $\Gamma_{\nu_0, \nu_1} = \mathbb{Z}/2\mathbb{Z}
\times \mathfrak{A}_5$. On a

$\alpha_0 = -1$, $\alpha_1$ \add{${} = \alpha$} solution de $\alpha^2 + \alpha
- 1 = 0$, donc

\struck{$\alpha_0 + \alpha_1$} $\alpha_0 + 2 = 1$, $\alpha_1 + 2$ \add{${} =
\alpha + 2$} inversible ($N(\alpha + 2) = -1 - 2 + 4 = 1$,
$1/(\alpha + 2) = \alpha' + 2 = -\alpha + 1$) $\alpha_0 + \alpha_1 = \alpha -
1$ inversible

[$N(\alpha - 1) = -1 + 1 + 1 = 1$, $1/(\alpha - 1) = \alpha' - 1 = -\alpha -
2$]. Donc
\note{il note le groupe alterné d'un A cerclé, rendu $\mathfrak{A}_5$ ;
$\alpha'$ est le conjugué de $\alpha$, racine de $\alpha^2 + \alpha - 1 = 0$}

\emph{Thm} Pour \ill{} \struck{\ill{}} couple $(G, \varphi)$, $G$ forme de
$\mathrm{GP}(1)_S$, $\varphi$ hom.\ $\Gamma = \mathbb{Z}/2 \times \mathfrak{A}_5
\to G$, soit $\alpha(\varphi)$ l'inv.\ de $\varphi(\rho_1)$. On trouve ainsi
une équivalence entre la cat.\ des $(G, \varphi)$ telles que $\forall s \in
S$, \add{$\varphi_s \mid \mathfrak{A}_5$ ne soit pas triviale},
\struck{$\rho_{0s}$ et $\rho_{1s}$ soient d'ordre} \struck{$\rho_0$ et
$\rho_1$ ne commutent pas} et la cat.\ discrète \struck{formée des}
définie par l'ens.\ des $\alpha \in \Gamma(S, \underline{\mathcal{O}}_S)$
satisfaisant
\[
\alpha^2 + \alpha - 1 = 0
\]
Les représentations $\varphi$ sont triviales sur $\mathbb{Z}/2$, donc
correspondent aussi aux
\[
\varphi^+ : \mathfrak{A}_5 \longrightarrow G
\]
qui ne sont triviales sur aucune fibre.

\marginal{Laisser tomber $\mathbb{Z}/2$ dans l'énoncé [celui-ci \ill{}
$\Leftrightarrow$ \struck{\ill{}} \ill{} \ill{} \ill{}]}
\note{note oblique dans la marge gauche, à la hauteur du théorème, reliée à
lui par une accolade ; sa seconde moitié est illisible}

\section*{Représentations projectives de $G_1$ (et repr.\ projectives
dégénérées de $G_2$)}

\note{inscrit en haut à droite du feuillet de garde (page 11), qui ne porte
rien d'autre et ne reçoit pas de numéro de page ; « dégénérées » est
souligné. À gauche, au crayon, « (10 p) » : le nombre de pages de la liasse
qui suit, pages 12 à 21}

\page{12}

\struck{\ill{}} \add{\emph{Prop.}} Soit $G$ forme de $\mathrm{GP}(1)$
\add{${} = \underline{\mathrm{Aut}}(X)$ ($X$ droite projective)} sur un corps
$k$ \add{alg.\ clos}, $\tau_0, \tau_1, \tau_2 \in G(k)$ tels que $\tau_0^2 =
\tau_1^2 = \tau_2^2 = (\tau_0\tau_2)^2 = 1$. Conditions équivalentes

A) $\exists$ \struck{diviseur} diviseur de degré 2 sur $X$ \struck{tel que}
\add{stable} par les $\tau_i$.

B) On peut trouver un sous-groupe \struck{\ill{}} $H$ de $G$ qui soit
\quad a) ou $N(T)$, $T$ tore maximal
\quad b) ou $B$, ss-groupe de Borel
tel que les $\tau_i \in H$

A') Les \struck{repr.\ projectives} prolongements \struck{projectifs}
\add{$\tau_i^C$ au plan projectif enveloppe $C$} \struck{\ill{}} ($i = 0, 1,
2$) de $\tau_i$ laissent une droite invariante ($\Leftrightarrow$ si car.\
${}\neq 2$ : laissent un pt invariant)
\note{des flèches courbes relient A) à B) et A') à C) dans la marge gauche}

\struck{C) L'image \struck{des} $G_2^+$ dans $G(k)$ [i.e.\ le groupe engendré
\ill{} $\rho_0 = \tau_0\tau_1$ et $\rho_1 = \tau_1\tau_2$, \ill{} par $\rho_0$
et $\sigma = \tau_0\tau_2 = \rho_0\rho_1$] est commutatif}
\note{C) est biffé et barré de traits obliques ; dans la marge gauche, dans
deux cadres reliés à lui, ses remplaçants :}
\marginal{C'') $\varphi$ se factorise par un groupe $\Gamma$ ext.\
\struck{d'un groupe} comm.\ \struck{\ill{}} par des groupes \emph{cycliques}}
\marginal{C') $\varphi : G_2 \to G(k)$ se factorise à travers un
$\mathbb{D}_\infty \times \mathbb{Z}/2$}

\struck{C)} D) Deux parmi les trois couples $(\tau_0, \tau_1)$ $(\tau_1,
\tau_2)$ $(\rho_0, \rho_1)$ commutent, i.e.\ \struck{\ill{}} $\rho_0^2 = 1$ ou
$\rho_1^2 = 1$ ou $\rho_0\rho_1 = \rho_1\rho_0$

\struck{D) Le \ill{} \ill{} \ill{} de $G(k)$ engendré par les $\tau_i$ est} \ill{}
isom., un $\mathbb{D}_n$ ($n \geqslant 1$ \uncertain{fini}, ou infini) ou
$\mathbb{D}_\infty \times \mathbb{Z}/2$ (\uncertain{impair})
\marginal{\emph{NB} dans (D), \ill{} \ill{} \ill{} \ill{} $\neq 2$, on peut
\ill{} $\mathbb{D}_n \times \mathbb{Z}/2$}
\note{dans la marge gauche, verticalement, encadré}

$A \Leftrightarrow A'$ : traduction via dictionnaire entre $X$ et $C$ (div.\
de degré 2 sur $X$ $\leftrightarrow$ droites de $C$)

$A \Leftrightarrow B$ via les stabilisateurs des div.\ de degré 2 sur $X$
$\mapsto$ les ss-groupes explicités dans B

On utilise que les groupes $H(k)$ sont des produits \uncertain{semi-directs}
$\gamma \cdot U$, $U$ comm., $\gamma$ \struck{\ill{}} \add{comm.}\
\struck{\ill{}} $_2\gamma$ d'indice 1 ou 2, \add{$\frac{1}{2}$} \ill{}
$_2U$ dans le centre de $\ldots$), \struck{\ill{}} \ill{} $\Gamma_1({}_2\gamma)
\cdot U$ (i.e.\ $_2\gamma = \lbrace 1 \rbrace$ ou $_2\gamma = \lbrace 1,
\varepsilon \rbrace$, \ill{} \ill{} l'identité sur $_2U$).
\note{« $_2\gamma$ » : la sous-partie de $\gamma$ tuée par 2, l'indice 2 est
écrit à gauche en bas ; « $\Gamma_1$ » est une lecture incertaine}

$B \Rightarrow C$ En effet, \struck{\ill{}} \struck{\ill{}} les images des
$\tau_i$ dans $\gamma$ \struck{\ill{}} \struck{\ill{}} si $_2\gamma = \lbrace 1
\rbrace$, \struck{\ill{}} \ill{} $H = B$, \ill{} car.\ 2] alors l'image
$\Gamma'$ de $\Gamma$ dans $G(k)$ \ill{} $\subset H(k)$ \ill{} est comm.,
\ill{} C. \ill{}, soit $\gamma = \lbrace 1, \varepsilon \rbrace$. Si
\struck{$(\tau_0, \tau_1)$ et $(\tau_1, \tau_2)$} \struck{$\tau_0$ et
$\tau_1$, $\tau_2$} \ill{} \struck{\ill{}, i.e.\ commutant, \ill{} \ill{}
$\tau_1$ $\tau_2$} \struck{\ill{} \ill{} on sait que \ill{}} \ldots\ on
aurait $\tau_i \in H(k)$, comme $\tau_i^2 = 1$, \ldots
\note{bas de page lourdement raturé ; la suite est à la page 13}

\page{13}

conduirait que $\tau_i$ est dans le centre --- ce qui \add{\uncertain{impliquerait}}
n'est pas le cas pour $i = 0$ ou 1 ou 2, si $(\tau_0, \tau_1)$ et $(\tau_1,
\tau_2)$ ne commutent pas. Dans ce cas, \ill{} \ill{} donc $\dot\tau_0,
\dot\tau_1, \dot\tau_2$ sont égaux : l'unique élément $\neq 1$ dans $\gamma$,
donc $\rho_0$ et $\rho_1$ sont dans $H$, donc commutent.
\note{les points sur $\dot\tau_i$ marquent les images dans $\gamma$}

$C \Rightarrow$ \struck{A} \add{C'}) C implique que l'image par
\uncertain{conjugaison} par les $\tau_i$ est isom.\ : un quotient d'un des
trois groupes dérivés de $G_2$ un divisant par l'une des relations
\struck{$(\tau_0\tau_1)^2 = 1$} $\rho_0^2 = 1$ \struck{$\tau_1\tau_2$}
$\rho_1^2 = 1$ ou $\rho_0\rho_1 = \rho_1\rho_0$. On voit \uncertain{tout} de
suite que chacun de ces groupes est isomorphe : $\Gamma' = \mathbb{Z}/2 \times
\mathbb{D}_\infty$ ; \struck{\ill{} \ill{}}

$C' \Rightarrow C''$\quad \struck{\ill{} \ill{} groupe, \ill{} \ill{}} un
ss-groupe invariant \add{cyclique} $\mathbb{D}_\infty^+ = \Gamma'' \subset
\Gamma'$ ; à quotient abélien

$C'') \Rightarrow A)$\quad Si $u$ est un générateur de $\Gamma''$, et si
$u_\varphi = \mathrm{id}$, l'hom.\ se factorise par un groupe comm.\
\emph{fini}, et on gagne (\ldots). Sinon, \struck{le} diviseur des pts fixes
de $u_\varphi$ est \uncertain{invariant} \uncertain{stable} \add{\uncertain{sous}}
\ill{} \ill{} de \struck{\ill{}} $\mathrm{Aut}(X)$ \uncertain{normalisant}
$u_\varphi$, donc stable par $\Gamma'$, on a ce qu'on \uncertain{voulait}.
\marginal{réf ?}

Donc A) $\Leftrightarrow$ C'') sont équivalentes d'ailleurs.
$D \Rightarrow C''$), montrons que B) \& C') impliquent D.\ Il suffit de
prouver que \add{$H$ ss-groupe de} \struck{$B$ ou $N(T)$} de $H = B$ ou
$N(T)$, qui \ill{} est quotient de $\mathbb{D}_\infty \times \mathbb{Z}/2$,
est isom.\ : $\mathbb{D}_\infty \times \mathbb{Z}/2$ \ldots\ ou : un
$\mathbb{D}_n$, ou : un $\mathbb{D}_n \times \mathbb{Z}/2$.
\note{la dernière phrase est d'une lecture incertaine dans son articulation :
les étiquettes $D \Rightarrow C''$ et « B) \& C') » sont écrites telles
qu'elles se présentent}

\page{14}

\emph{Lemme} Soit $\mathbb{D}_\infty = \lbrace \sigma_0, \sigma_1 \mid
\sigma_0^2 = \sigma_1^2 = 1 \rbrace$, $u = \sigma_0\sigma_1$ \add{groupe
diédral infini}, \struck{et soient les ss-groupes $N$ de $\mathbb{D}_\infty$,
\ill{} \ill{}, sont les suivants \ill{}} \struck{\ill{}} $i \in
\mathbb{N}^*$ (entier ${}\geqslant 0$)
\note{le début de l'énoncé est très raturé ; plusieurs rédactions
successives sont biffées}

\struck{a) Soit \ill{} $N_i$ $\forall i \in \mathbb{N}^*$ (\ill{})}, soit
$N_i$ le sous-groupe engendré par $u^i$, c'est un ss-groupe invariant,
\add{cyclique infini ($i \neq 0$) ou 1}, \ill{} quotient
$\to \mathbb{D}_\infty/N_i$ \ldots\ \uncertain{Donc} on a :
\[
\mathbb{D}_i = \lbrace \sigma_0, \sigma_1 \mid \sigma_0^2 = \sigma_1^2 =
(\sigma_0\sigma_1)^i = 1 \rbrace
\qquad (\emph{NB}\ \text{la déf.}\ \mathbb{D}_0 = \mathbb{D}_\infty)
\]

b) \struck{Soit \ill{} $\forall i \in \mathbb{N}^*$, soit $N'_i$ de
$\mathbb{D}_\infty$ engendré par $u^i$ et} Soit $i \in \mathbb{N}^*$, et soit
$j \in \mathbb{Z}/i\mathbb{Z}$. \struck{Considérons} Notons $N_{i,j}$ le
sous-groupe de $\mathbb{D}_\infty$ engendré par $u^i$ et $\sigma_0 u^{j'}$,
avec $j' \in \mathbb{Z}$ ($j'$ mod $i$) ${}= j$ ; c'est un sous-groupe
\add{isom.\ à $\mathbb{D}_\infty$ ($i \neq 0$), ou $\mathbb{Z}/2\mathbb{Z}$
($i = 0$), \ill{}} invariant ssi $\mathbb{D}_i$ est commutatif, i.e.\ $i \in
\lbrace 1, 2 \rbrace$ \struck{\ill{}}
\struck{\ill{} conjugué de $\sigma_0$ ssi $j \in 2\mathbb{Z}/i\mathbb{Z}$}
\struck{(p.ex.\ si $i$ impair) un conjugué de $\sigma_1$}
[\uncertain{De} plus, si $i$ est pair, \add{\ill{}} un groupe contient des
conjugués de $\sigma_0$ et $\sigma_1$ (\struck{qui sont \ill{}} des images
sous conjugaison)].
\struck{Je \ill{} \ill{}} Quel que soit $i$, $N_{i,j}$ contient un conjugué
de $\sigma_0$ ssi $j \in 2\mathbb{Z}/i\mathbb{Z}$, un conjugué de $\sigma_1$
ssi $(j + 1) \in 2\mathbb{Z}/i\mathbb{Z}$ \struck{\ill{}}

c) Les ss-groupes $\lbrace N_i, N_{i,j} \rbrace$ sont tous \ill{} deux à deux
\uncertain{distincts} et ce sont les ss-groupes de $\mathbb{D}_\infty$
\uncertain{sauf} \ill{} \uncertain{un} \ill{} \ill{}.

\emph{P.\ ex.} $\lbrace 1 \rbrace = N_0$, $\mathbb{D}_\infty = N_{1,0}$

\emph{Cor 1} Les \add{types d'isom.\ des} groupes quotients de
$\mathbb{D}_\infty$ sont \add{\uncertain{ceux}} des $\mathbb{D}_i$ ($i
\geqslant 1$, \uncertain{ou} infini), \struck{\ill{}} et \struck{\ill{}} \ldots\
des gpes \uncertain{unité}.

\emph{Cor 2} \struck{\ill{} \ill{}} \add{\uncertain{Types}} d'iso.\ des
groupes quotients \struck{de $\mathbb{D}_\infty \times \mathbb{Z}/2$
\ill{} \ill{} \ill{} \ill{}} de $\mathbb{D}_\infty \times \mathbb{Z}/2$ ---
les $\mathbb{D}_i \times \mathbb{Z}/2$ (et les $\mathbb{D}_i$ \ill{}),
\struck{\ill{} \ill{} faisant $\mathbb{D}_\infty$ \ill{} \ill{}}
\marginal{\emph{NB} $\mathbb{D}_i \times \mathbb{Z}/2 \simeq \mathbb{D}_{2i}$
si $i$ impair}

\struck{1--1 \ill{} \ill{} \ill{} \ill{} $N \subset N'$ de deux ss-groupes
de $\mathbb{D}_\infty$ \ill{} \ill{} $\mathbb{D}_\infty$ \ill{} \ill{}
$(N = M \cap \mathbb{D}_\infty$, $N' = \mathrm{Im}(M \to
\mathbb{D}_\infty))_{\ldots}$ \ill{}}
\note{encadré et barré de traits obliques : une première tentative de décrire
les sous-groupes $M$ de $\mathbb{D}_\infty \times \mathbb{Z}/2$ par un couple
$N \subset N'$ de sous-groupes de $\mathbb{D}_\infty$ ; la ligne qui le suit
(« \ldots\ invariants ssi $N$, $N'$ le sont, \ldots ») est biffée aussi}

\struck{\ill{}} \emph{Dém} C'est \uncertain{évident} si des \uncertain{produits}
de $\mathbb{D}_\infty \times \mathbb{Z}/2$ (par $N_i$ ou $N_i \times
\mathbb{Z}/2$). Si $\Gamma = (\mathbb{D}_\infty \times \mathbb{Z}/2)/N$ est un
groupe quotient, \ill{} $N$ ss-groupe \ill{}, si $N \subset \mathbb{D}_\infty$
\ill{} \ill{}. \ill{} \ill{} $\Gamma \simeq \mathbb{D}_\infty/N \times
\mathbb{Z}/2$ et \ldots\ implique \uncertain{cas} 1, \uncertain{si} non
$N \xrightarrow{\ \mathrm{pr}_2\ } \mathbb{Z}/2$ \ldots

\page{15}

donc $(\mathbb{D}_\infty \times \mathbb{Z}/2)/N \simeq \mathbb{D}_\infty/N
\cap \mathbb{D}_\infty$ et \ldots\ \uncertain{implique} \ill{} comme le
\uncertain{corollaire}.

\emph{NB} Les groupes \struck{$\mathbb{D}_i \times \mathbb{Z}/2$},
$\mathbb{D}_i \times \mathbb{Z}/2$ ($i$ infini ou pas pair) et $\mathbb{D}_i$
sont ils \uncertain{effectivement} \uncertain{plongeables} dans
\add{(\uncertain{infini}) gp.\ diédraux} un $H = N(T)$ ou $B$ ?
\struck{Parmi les $\mathbb{D}_i$, \ill{}} \struck{diédraux pour \ill{}
$N(T)$, \uncertain{pourvu} que le seul \ill{} gp.\ des} \struck{\ill{} \ill{}
\ill{} mais le cas pour $\mathbb{D}_i \times \mathbb{Z}/2$ \ill{} de car.\
\uncertain{première} à 2 \ldots} \struck{\ill{} $\mathbb{D}_2 \simeq
\mathbb{Z}/2\mathbb{Z} \times \mathbb{Z}/2\mathbb{Z}$ \ill{} \ill{} \ill{}}
\add{$\simeq V_4$} \struck{\ill{} la \uncertain{rep.}\ diédrale.} On
\uncertain{considère} le \ldots
\marginal{explicitez \uncertain{ce} \ill{} des représentations \uncertain{non}
diédrales}
\note{note oblique dans la marge gauche, reliée par une accolade à ce
\emph{NB} ; tout le paragraphe est coupé de ratures, et la ligne suivante
(« \emph{Proposition} \ldots\ que les $\mathbb{D}_i$ admettent une repr.\
\ill{} dans $\mathrm{GP}(1, k)$, il faut que \ldots ») est barrée de traits
obliques et biffée}

\emph{Proposition} Soit $k$ corps alg.\ clos (de car.\ $p$), $i \in
\mathbb{N}^*$, \struck{\ill{}} conditions \ill{} \ill{} conditions

a) $\exists$ repr.\ fidèle $\mathbb{D}_i \hookrightarrow \mathrm{GP}(1, k)$
\note{« $\mathbb{D}_i$ » est surchargé les deux fois}

\emph{Alors} : \add{\uncertain{ou} $i = 1$}

a) si $p = 0$, (a) est satisfait. \struck{\ill{}}

\struck{b) si $p \neq 0$, (i.e.\ $\mathbb{D}_i = \mathbb{D}_0 = \mathbb{D}_\infty$)}
\add{Supposons $p > 0$, si $i = 0$ (i.e.\ $\mathbb{D}_i = \mathbb{D}_0 =
\mathbb{D}_\infty$)} (a) équivaut à la condition : \struck{\ill{}} $k$
\uncertain{non} algébrique sur le corps premier (\uncertain{ou encore} :
$k^*$ contient un élément d'ordre infini)
\note{la parenthèse « Supposons \ldots\ $\mathbb{D}_\infty$) » est cerclée et
remplace le b) biffé}

\struck{b)} \add{\uncertain{c)} si $p \neq 0$, et}

b) si $i \geqslant 2$, (a) équivaut à la condition :

\struck{\ill{}} (**) On a \struck{$i$ premier} $p \nmid i$, ou bien $i = p$,
\struck{\ill{}}

\struck{ou bien $i = 2p$ et $p$ impair.}

\emph{Dém.} a), b) immédiats en \struck{cas} que $i = 0$ car
équivalent à l'existence d'un él.\ d'ordre infini dans $k^*$ [ou dans
$k$].
\struck{Pour $i \geqslant 2$, on sait que les} \struck{d'ordre fini}
\struck{c) On sait que les ordres des él.\ fts (de $\mathrm{GP}(1)$ sont
\ill{} les \ill{} \ill{} premiers, \ill{} \ill{} \ill{}} \struck{b) les
\ill{} \ill{} \ill{}, \ill{} $H = \gamma \cdot U$ (\ill{} \ill{})}
\marginal{explicitez \uncertain{ceux}}
\note{le bas de la page est biffé ligne à ligne et barré de traits
obliques ; en marge gauche, « explicitez » et un mot incertain}

\page{16}

\struck{Car c) si $H(k) = \gamma(k) \cdot U(k)$ ($U(k) \simeq k$ \ldots)}

$*) \Rightarrow (**)$ \struck{si $p$ impair}, car les ordres des éléments
d'ordre fini de $\mathrm{GP}(1)$ sont les $i$ \uncertain{premiers} à $p$
\struck{et $p = 2$} (correspondant aux él.\ \uncertain{semi-simples}) et $p$
(correspondant aux élts unip.\ ${}\neq 1$) ; \struck{Mais si $i = p = 2$,
\ill{} $\mathbb{D}_i = \mathbb{Z}/2 \times \mathbb{Z}/2$} et \ldots\ les
\ill{} sont contenus dans des \ldots\ de $\mathrm{GP}(1)$ (car \ill{} les
$G_m$ \uncertain{sont} les \uncertain{tores} \uncertain{maximaux}). Et

$(**) \Rightarrow (*)$ : immédiat, en \uncertain{distinguant} et
\uncertain{examinant} (pour le cas $i = p$) \uncertain{le} cas $p$ impair, et
le cas \uncertain{particulier} où $\mathbb{D}_2$ \uncertain{est} \struck{\ill{}}
\ill{} dans $U$ !, \uncertain{voir} \uncertain{Remarque}.
\note{la fin de cette ligne est serrée entre deux lignes et peu lisible}

\struck{\emph{Cor} Soit $i$ pair, \ldots\ Supposons (\ldots)
(*') $\exists$ repr.\ fidèle $\mathbb{D}_i \times \mathbb{Z}/2
\to \mathrm{GP}(1, k)$.
a) Si $p = 0$, \struck{\ldots\ $i = 1$} (*') est satisfait ssi $i = 0$
\struck{\ill{}} \add{\uncertain{donc} $\mathbb{D}_i \times \mathbb{Z}/2 \simeq
\mathbb{D}_\infty \times \mathbb{Z}/2$}\quad
b) Si $i = 0$, (*') est satisfait ssi \struck{\ill{}} $p \neq 2$.\quad
c) Si $i \geqslant 2$, $p > 0$, alors (*') n'est pas satisfait.}
\note{ce corollaire est barré de deux grandes croix et encadré d'un trait à
gauche}

\emph{Dém} \struck{\ill{} \ill{}} \struck{Considérons}
\[
G = G_1 = \lbrace \tau_0, \tau_1 \mid \tau_0^2 = \tau_1^2 = 1 \rbrace
\]
On s'intéresse aux repr.\ de $\mathbb{D}_\infty$ \add{$= G_1$} dans des $G$
formes de $\mathrm{GP}(1)$ (\add{$= \underline{\mathrm{Aut}}(X)$}) (\uncertain{sur}
$S$) $\Leftrightarrow$ la cat.\ des collections \uncertain{équivalentes}
\marginal{on \uncertain{pose} $t_i = \varphi(\tau_i)$ ($i = 0, 1$), $u =
\tau_0\tau_1$ (${}= \varphi(u)$)}
: \uncertain{celle} des représentations \add{\ill{}} projectives de dim.\ 2
(\add{dans un plan projectif $X$}) de $G_1$, \uncertain{munies} de
\struck{\ill{}} \ill{} quelconques \uncertain{bien} \uncertain{connues}. On
va regarder d'abord les cas où la condition \struck{On \ill{} \ill{}}
\struck{\ill{} \ill{}} suivante est satisfaite

\struck{a) $\forall s \in S$, $u_s$ ($= \tau_{0s}\tau_{1s}$) \ill{}}

a) $\forall s \in S$, $u_s^2 \neq 1$ (i.e.\ $\varphi_s : \mathbb{D}_\infty \to
G_s$ ne se factorise pas par \uncertain{un} $\mathbb{D}_2$)

i.e.\ $\tau_{0s}$, $\tau_{1s}$ ne commutent pas

\struck{La catégorie des \ill{} \ill{} \ill{} couples $(G, \varphi)$
\uncertain{équivaut}} \ill{} \ill{}, \ill{} \ill{} \ill{} introduit $a_0$,
$H_0$, $a_1$, $H_1$, comme \struck{à la cat.\ des polygones réguliers}
d'habitude, la condition a) implique
\[
a_{0s} \notin H_1 \ (\Leftrightarrow a_1 \notin H_0) \ \text{sur chaque fibre}
\]
\ldots\ quand \ldots
\note{la page s'arrête sur une ligne inachevée, reprise à la page 17}

\page{17}

dans \add{\uncertain{\ill{}}} loc.\ (cas particulier d'une situation d'un
\emph{polygône régulier} \struck{\uncertain{projectif}}) dans le \ill{}
projectif $X/S$, tel que \struck{\ill{}} $\beta = \alpha + 2$ \uncertain{inversible}
(ceci exprime $a_1 \notin H_0$), on \uncertain{\ill{}} \struck{\ill{}} \ill{}
le quadrique \ill{} \uncertain{associée} \add{$Q$} \uncertain{dans} $X$
(\uncertain{l'unique} quadrique) soit lisse. Donc
\note{« $\beta = \alpha + 2$ inversible » est encadré ; « polygône régulier »
est souligné. La lecture de la ligne sur la quadrique est incertaine}

\struck{\ill{} \ill{} \ill{} \ill{} \ill{} \ill{} \ill{}}
\note{une ligne entièrement raturée}

\emph{Th} Soit \struck{\ill{}} $C$ la catégorie des couples $(G, \varphi)$
d'une forme $G$ de $\mathrm{GP}(1)_S$ et d'un hom.\ $\varphi : \mathbb{D}_\infty
\to G(S)$ [\uncertain{d'où} $\tau_0, \tau_1, u \in G(S)$] satisfaisant

a) $\forall s \in S$, $u_s^2 \neq 1$ (i.e.\ $\varphi_s$ ne se factorise pas par
$\mathbb{D}_2$)

b) \struck{Soit $C$ la catégorie des couples formés d'un \ill{}} $\forall \rho
\in \Gamma(S, \mathcal{O}_S^*)$, soit $U(\rho)$ le groupe \struck{\ill{}}
lisse de dim.\ 1 sur $S$, \struck{défini} formé des $\lambda$ tels que $1 +
\lambda\rho$ inversible, avec la multiplication
\[
[\lambda] \cdot [\lambda'] = [\lambda + \lambda' + \lambda\lambda'\rho]
\]
Soit $C'$ la catégorie des couples $(\alpha, R)$, où $\alpha \in \Gamma(S,
\mathcal{O}_S)$ telle que $\beta = \alpha + 2$ inv., et $R$ un torseur sous
\struck{\ill{}} $U(2 - \alpha)$.
\marginal{\ill{} \uncertain{Sinon} \ill{} \ill{} \ill{} \ill{} \ill{}
$\mathbb{D}_\infty$ \ill{} \ill{} \ill{} \ill{} \ill{} $\times \mathbb{D}(G)$
\ill{} \ill{}}
\note{note oblique sur quatre lignes dans la marge gauche, face à b) ; très
peu lisible}

c) Soit $\alpha \in \Gamma(S, \mathcal{O}_S)$, \struck{\ill{} de $\mathcal{O}_S$}
et considérons la \ill{} \add{si $\alpha + 2$ inv.\ dans $\mathcal{O}_S$ lisse,}
représentation \uncertain{privilégiée} de \uncertain{invariant} $\alpha$,
\ill{} $(\ldots)$, \ill{} \ill{} \ill{} \ill{} : $G_\alpha \simeq
\mathrm{Aut}(X_\alpha, Q_\alpha)$ \ill{},
$\varphi_\alpha : \mathbb{D}_\infty \to G_\alpha(S)$ \struck{\ill{} \ill{}
$U_\alpha$ opère $\to \mathrm{Aut}(G_\alpha, \varphi_\alpha)$}.
\[
U_{2-\alpha} \xrightarrow{\ \sim\ } \mathfrak{z}\,\mathrm{Ant}(\varphi_\alpha) =
\underline{\mathrm{Aut}}(G_\alpha, \varphi_\alpha)
\]
Si $R$ est un torseur sous $G_\alpha$, on peut l'utiliser pour former
$(G_\alpha^R, \varphi_\alpha^R)$.
\note{la ligne « $U_{2-\alpha} \simeq \ldots$ » est ajoutée au-dessus de la
ligne biffée ; « $\mathfrak{z}\,\mathrm{Ant}$ » (centralisateur ?) est une
lecture incertaine}

d) On trouve ainsi un foncteur
\[
C' \longrightarrow C,
\]
et celui-ci est une équivalence de cat.\ Le foncteur quasi-inverse est
\ldots

\page{18}

le foncteur $C \to C'$ défini ainsi : donné $(G, \varphi) \in \mathrm{Ob}\,C$,
on définit $\alpha$ comme l'invariant de $\varphi(\underline{u}) = u$, et $R$
comme le \struck{formé des \ill{}} faisceau de $X$ (\add{\uncertain{droite
projective}} \ill{} de $C$ dont \uncertain{proj.}\ $X/S$ défini par $G$)
\struck{formé de} $H_1 \setminus \lbrace c \cup d \rbrace$, qui est un torseur
sous $U_{2-\alpha} \xrightarrow[\text{iso can}]{\ \sim\ }
\underline{\mathrm{Aut}}(X, \varphi)$ \ldots\ tenant :
$\underline{\mathrm{Aut}}(Q_X, \varphi)$, c'est le ss-groupe
\struck{\uncertain{Corollaire} Soit \ill{} \ill{} La \uncertain{relation}}
\note{« $\underline{\mathrm{Aut}}(X, \varphi)$ » surcharge une première
écriture raturée}

\struck{de $U_\rho$ formé des $[\lambda]$ tels que $\lambda
\mathbin{\underset{\rho}{*}} \lambda = 1$. Donc $\lambda(2 + \lambda\rho) = 0$
i.e.\ \ldots\ (\uncertain{\ill{}}) $U_\rho \simeq \mathbb{G}_m$ (si $\alpha
\neq 2$) et $\rho \neq 2$. Alors}
\note{ce passage est encadré à gauche d'un trait et fermé en bas par une
ligne courbe ; il semble abandonné}

1) Cas $2\rho \neq 0$ \add{\uncertain{inversible}} \struck{i.e.\ $\alpha \neq
2$}, \struck{i.e.} (\struck{\ill{} \ill{} \ill{}. Alors}
$\underline{\mathrm{Aut}}(X, \varphi) \simeq$ \struck{\ill{}} $\subset
\mathbb{G}_m$, \ill{} \uncertain{engendré} par $\sigma$. D'autre part l'image
de $U_\alpha$ dans \struck{\ill{}} $G$ \struck{\ill{}} est l'unique
\uncertain{tore} maximal qui contient \struck{l'image} $u$ ; dans
\ill{} si $u$ a d'ailleurs pour \uncertain{partie} $2 \ldots$ \ldots\
$\sigma = u^{m}$.
\note{deux grands traits obliques barrent ce 1) et la suite, peut-être
pour l'annuler}

2) Cas $\rho = 0$, \struck{$\rho \neq 2$} \struck{$2 \neq 0$} (i.e.\ $U_\rho
\simeq \mathbb{G}_a$ i.e.\ $\alpha = 2$) et \struck{$p = 2$} 2 inv.
[\struck{Alors, l'image de} $U_\rho$ dans $G$ est l'unique
\struck{\ill{}} \ill{} $\simeq \mathbb{G}_a$ qui contient
\struck{\ill{} \ill{}} \ldots\ $\underline{\mathrm{Aut}}(X, \varphi) =$ \ill{}
\ill{} \ldots\ $u$]
\note{le bas de la page est raturé en plusieurs couches}

\marginal{\ill{} des \ill{} \ill{} d'un \ill{} \ill{} \ill{} $\mathbb{D}_\infty
\times \mathbb{Z}/2$ : \ill{} \ill{} représentations \ill{} \ill{}
$\mathbb{D}_i$ \ill{} \ill{} \ill{} diédrales}
\note{deux notes obliques dans la marge gauche, entre des traits verticaux ;
lecture très partielle}

\page{19}

\struck{3) $\rho \neq 0$, $\rho = 2$} \struck{\add{inv., \uncertain{soit} $\alpha \neq
2$}} \struck{(i.e.\ $U_\rho \simeq \mathbb{G}_m$) \ill{}}
\struck{$\mathrm{Aut}(X, \varphi) = \lbrace 1 \rbrace$}
\struck{$\underline{\mathrm{Aut}}(Q, \varphi) \simeq {}_2U_\rho \simeq \mu_2$}
\struck{On conclut}
\note{ce 3) est barré de traits obliques ; la dernière égalité est en partie
effacée}

$\underline{\mathrm{Aut}}(G, \varphi) \simeq {}_2(U_\rho)$

\struck{\emph{Proposition}}

\emph{Corollaire 1} Supposons $\rho$ \add{${}= 2 - \alpha$} inv., i.e.\
$U_\rho \overset{\text{loc}}{\simeq} \mathbb{G}_m$, i.e.\ \struck{a) si
$\rho$ inv., et $2 \cdot 1_S$ \ill{} \ill{}, alors}
\[
U_\rho \xrightarrow[\text{can}]{\ \sim\ } \mathbb{G}_m, \qquad \lambda
\longmapsto 1 + \lambda\rho .
\]
Alors : a) l'image $T$ de $U_\rho$ dans $G$ est l'unique tore maximal de $G$
contenant $u$ ;

\struck{b) si $2 \mid \rho$ \ill{} \ill{} $2 \mid \rho = 0$, alors}
b) $\underline{\mathrm{Aut}}(G, \varphi) = {}_2T \overset{\text{can}}{\simeq}
\mu_2$
\note{« $\overset{\text{loc}}{\simeq}$ », « can » : il écrit « loc » et
« can » au-dessus du signe}

\emph{Corollaire 2} Supposons $\rho = 2 - \alpha = 0$ (i.e.\ $\alpha = 2$),
i.e.\ \struck{$U_\rho \simeq$} $U_\rho \xrightarrow{\ \sim\ } \mathbb{G}_a$.
Alors a) l'image $U$ de $U_\rho$ \struck{\ill{}} est un
\uncertain{sous-groupe} \struck{maximal \ill{}} unipotent maximal de $G$
\uncertain{contenant} \ill{}

b) Si $2 \cdot 1_S$ \uncertain{inv.}, alors $\underline{\mathrm{Aut}}(G,
\varphi) = \lbrace 1 \rbrace$

c) Si $2 \cdot 1_S = 0$ i.e.\ $S$ de car.\ 2, alors
\[
\underline{\mathrm{Aut}}(G, \varphi) = U \simeq \mathbb{G}_a .
\]

\emph{Cor.\ 3} Soit $\gamma$ un groupe, et soit $\Gamma = \mathbb{D}_\infty
\times \gamma$. Considérons la cat.\ $C(\gamma)$ des couples $(G, \varphi)$, $G$
forme de $\mathrm{GP}(1)_S$, $\varphi : \Gamma \to G(S)$. Cette catégorie est
équivalente : elle \ldots
\note{la page s'arrête sur cette phrase inachevée}

\page{20}

des triples $(\alpha, R, \varphi_0)$, où a) $\alpha \in \Gamma(S,
\mathcal{O}_S)$ tel que $\alpha + 2$ inv.\ (i.e.\ $\alpha(s) \neq -2$ $\forall
s \in S$), b) $R$ un torseur sous \add{le groupe} $U_{\rho = \alpha - 2}$,
c) $\varphi_0 : \gamma \to ({}_2U_\rho)(S)$ un homomorphisme de groupes.
\note{« $U_{\rho = \alpha - 2}$ » est écrit tel quel ; les pages 17 et 19
posent $\rho = 2 - \alpha$}

\emph{Proposition} Considérons \uncertain{homomorphismes} \struck{\ill{}
\ill{}} \struck{\emph{Cor.\ Unicité} Considérons un $\Gamma = \mathbb{D}_i
\times \mathbb{Z}/2$ (\ill{})} et un corps alg.\ clos $k$, \uncertain{enfin}
la condition
\note{le titre « Proposition » est écrit au-dessus d'un premier titre biffé ;
les deux lignes sont encadrées à gauche}

(*) $\exists$ hom.\ injectif $\Gamma = \mathbb{D}_i \times \mathbb{Z}/2
\hookrightarrow \mathrm{GP}(1, k)$.

\struck{a) Si $i = 0$, alors (*) est satisfait \ill{}}

\ldots\ \uncertain{solution} $\varphi : \mathbb{D}_\infty \to G$ ($G$ forme de
$\mathrm{GP}(1)$).

\emph{Cond.} Soit $n \in \mathbb{N}^*$, \add{$n \geqslant 3$}, considérons
$\mathbb{D}_n \simeq \mathbb{D}_\infty/\lbrace u^n \rbrace$. Considérons les
conditions (équivalentes)

a) $u^n = 1$, et $\forall s \in S$, $u_s$ d'ordre $n$ \struck{[\ill{} $n =
0$, il faut dire : $u_s$ d'ordre $\infty$]}

b) $\varphi$ se factorise par un \emph{mono} $\mathbb{D}_n \hookrightarrow G$.

On va les exprimer en termes de l'invariant $\alpha$. On trouve :
\struck{\ill{} \ill{} \ill{}}

\struck{1) Si $n = \infty$ : il faut et il suffit que $\forall s \in S$,
\ill{} \ill{}} \struck{c) $\forall s \in S$ \ill{}}
\note{ce 1) est barré et encadré d'un trait sinueux}

c) 1) $n \cdot 1_S$ est inversible, \emph{ou} $n = p$ \uncertain{premier}
\struck{\ill{} \ill{} \ill{} \ill{}} \struck{\ill{} \ill{} \ill{}}
\uncertain{est} premier impair, \struck{\ill{}}) \emph{et} 2) $\Psi_n(\alpha) =
0$ \quad [$\Psi_n \in \mathbb{Z}[U]$ polynôme \uncertain{½cyclotomique}
d'ordre $n$]
\note{« $\Psi_n$ » : la même lettre qu'à la page 7. « ½cyclotomique » :
son abréviation de \emph{semi}-, lecture probable}

\end{document}
